% SEGMENT OLP-0632-S01
% Part: history
% Chapter: biographies
% Section: ernst-zermelo

\documentclass[../../../include/open-logic-section]{subfiles}

\begin{document}

\olfileid{his}{bio}{zer}

\olsection{எர்ன்ஸ்ட் செர்மெலோ}

\olphoto{zermelo-ernst}{எர்ன்ஸ்ட் செர்மெலோ}

எர்ன்ஸ்ட் செர்மெலோ 1871 ஜூலை 27 அன்று ஜெர்மனியின் பெர்லினில் பிறந்தார்.
அவருக்கு ஐந்து சகோதரிகள் இருந்தனர்; ஆனால் குடும்பத்தில் உடல்நலக்
குறைபாடுகள் இருந்ததால் அவர்களில் மூவர் மட்டுமே முதிர்வயதுவரை
உயிர்வாழ்ந்தனர். அவருடைய பெற்றோரும் இளம் வயதிலேயே இறந்தனர்;
செர்மெலோவுக்கு 17 வயதானபோது அவரும் உடன்பிறந்தவர்களும் பெற்றோரற்றவர்களாக
இருந்தனர். கலைகளில், குறிப்பாகக் கவிதையில், அவருக்கு ஆழ்ந்த ஆர்வம்
இருந்தது. கூர்மையான அறிவும் நகைச்சுவையும் விமர்சனப் பார்வையும்
உடையவர் என்று அறியப்பட்டார். 1904-இல் தேர்வு அடிகோளை அறிமுகப்படுத்தியதும்,
1908-இல் கணக்கோட்பாட்டை அடிகோள்மயமாக்கியதும் அவருடைய புகழ்பெற்ற
கணிதச் சாதனைகளில் அடங்கும்.

% SEGMENT OLP-0632-S02
பல்கலைக்கழகத்தில் செர்மெலோவின் ஆர்வங்கள் பலதரப்பட்டவை. இயற்பியல்,
கணிதம், மெய்யியல் ஆகியவற்றைப் படித்தார். ஹெர்மன் ஷ்வார்ட்ஸின்
மேற்பார்வையில் பெர்லின் பல்கலைக்கழகத்தில் \emph{Investigations in
  the Calculus of Variations} என்னும் ஆய்வுரையை 1894-இல் முடித்தார்.
1897-இல் க\"{o}ட்டிங்கன் பல்கலைக்கழகத்தில் மேல்படிப்பைத் தொடர முடிவு
செய்தார்; அங்கே டேவிட் ஹில்பர்ட்டின் அடித்தளவியல் பணியால் பெரிதும்
தாக்கம் பெற்றார். 1899-இல் பேராசிரியர் பதவிக்குத் தகுதி பெற்றாலும்,
பதினோரு ஆண்டுகளுக்குப் பிறகே அது கிடைத்தது; அவரது வழக்கத்திற்கு
மாறான நடத்தையும் ``பதற்றமான அவசரமும்'' இதற்குக் காரணமாக
இருந்திருக்கலாம்.

% SEGMENT OLP-0632-S03
1910-இல் சூரிக் பல்கலைக்கழகத்தில் ஊதியமுள்ள பேராசிரியர் பதவியை
இறுதியாகப் பெற்றார்; ஆனால் காசநோயால் 1916-இல் ஓய்வு பெற வேண்டியிருந்தது.
குணமடைந்தபின் 1921-இல் ஃப்ரைபர்க் பல்கலைக்கழகத்தில் கௌரவப்
பேராசிரியர் பதவி வழங்கப்பட்டது. அக்காலத்தில் கணிதத்தின் அடித்தளங்கள்
குறித்துப் பணியாற்றினார். தோரால்ஃப் ஸ்கோலெம், கர்ட் கோடல் ஆகியோரின்
பணிகள் அவருக்கு எரிச்சலூட்டின; தமது கட்டுரைகளில் அவர்களுடைய
அணுகுமுறைகளை வெளிப்படையாக விமர்சித்தார். அவரது செல்வாக்கின்மையாலும்
ஜெர்மனியில் ஹிட்லர் ஆட்சிக்கு வந்ததை அவர் எதிர்த்ததாலும்,
1935-இல் ஃப்ரைபர்க் பதவியிலிருந்து நீக்கப்பட்டார்.

% SEGMENT OLP-0632-S04
செர்மெலோவின் பிற்கால வாழ்க்கை தனிமையால் குறிக்கப்பட்டது. 1935-இல்
பதவி நீக்கப்பட்டபின் கணிதத்தை விட்டுவிட்டார். கிராமப்புறத்துக்குச்
சென்று எளிமையாக வாழ்ந்தார். 1944-இல் திருமணம் செய்துகொண்டார்;
பார்வை குறைந்துவந்ததால் மனைவியை முற்றிலும் சார்ந்திருக்க வேண்டிய
நிலை ஏற்பட்டது. 1951-க்குள் பார்வையை முழுவதும் இழந்தார். 1953
மே 21 அன்று ஜெர்மனியின் க\"{u}ன்டர்ஸ்டாலில் மறைந்தார்.

% SEGMENT OLP-0632-S05
\begin{reading}
செர்மெலோவின் விரிவான வாழ்க்கை வரலாற்றுக்கு \citet{Ebbinghaus2015}-ஐப்
பார்க்கவும். 1904, 1908 ஆண்டுகளில் வெளிவந்த அவருடைய முக்கியமான
கட்டுரைகளை மூல ஜெர்மன் மொழியில் படிக்கலாம்
\citep{Zermelo1904,Zermelo1908}. இயற்பியல் பற்றிய எழுத்துகள்
உட்படத் தொகுக்கப்பட்ட படைப்புகள் ஆங்கில மொழிபெயர்ப்பில் கிடைக்கின்றன
\citep{Ebbinghaus2010,Ebbinghaus2013}.
\end{reading}

\end{document}
